\begin{ex}
Góc vuông có số đo là bao nhiêu độ?
\shortans{90}
\loigiai{
$90^\circ$.
}
\end{ex}

\begin{ex}
Góc bẹt có số đo là bao nhiêu độ?
\shortans{180}
\loigiai{
$180^\circ$.
}
\end{ex}

\begin{ex}
Góc nhọn có số đo bé hơn bao nhiêu độ?
\shortans{90}
\loigiai{
Bé hơn $90^\circ$.
}
\end{ex}

\begin{ex}
Góc tù có số đo lớn hơn bao nhiêu độ và bé hơn $180^\circ$?
\shortans{90}
\loigiai{
Lớn hơn $90^\circ$.
}
\end{ex}

\begin{ex}
Lúc 3 giờ đúng, góc giữa kim giờ và kim phút bằng bao nhiêu độ?
\shortans{90}
\loigiai{
$90^\circ$.
}
\end{ex}

\begin{ex}
Lúc 6 giờ đúng, góc giữa hai kim bằng bao nhiêu độ?
\shortans{180}
\loigiai{
$180^\circ$.
}
\end{ex}

\begin{ex}
Lúc 9 giờ đúng, góc giữa hai kim bằng bao nhiêu độ?
\shortans{90}
\loigiai{
$90^\circ$.
}
\end{ex}

\begin{ex}
Góc bẹt bằng mấy lần góc vuông?
\shortans{2}
\loigiai{
$180 : 90 = 2$ lần.
}
\end{ex}

\begin{ex}
Góc $60^\circ$ là góc nhọn hay tù? (Nhọn ghi 1, tù ghi 2).
\shortans{1}
\loigiai{
Góc nhọn (1).
}
\end{ex}

\begin{ex}
Góc $135^\circ$ là góc nhọn hay tù? (Nhọn ghi 1, tù ghi 2).
\shortans{2}
\loigiai{
Góc tù (2).
}
\end{ex}

\begin{ex}
Một hình tam giác vuông có bao nhiêu góc vuông?
\shortans{1}
\loigiai{
Có 1 góc vuông.
}
\end{ex}

\begin{ex}
Một hình chữ nhật có bao nhiêu góc vuông?
\shortans{4}
\loigiai{
Có 4 góc vuông.
}
\end{ex}

\begin{ex}
Tổng số đo 3 góc trong một tam giác bằng bao nhiêu độ?
\shortans{180}
\loigiai{
Tổng 3 góc bằng $180^\circ$.
}
\end{ex}

\begin{ex}
Một tam giác có 1 góc vuông và 1 góc $40^\circ$. Góc còn lại có số đo bao nhiêu độ?
\shortans{50}
\loigiai{
$180 - 90 - 40 = 50^\circ$.
}
\end{ex}

\begin{ex}
Một tam giác đều có 3 góc bằng nhau. Mỗi góc có số đo bao nhiêu độ?
\shortans{60}
\loigiai{
$180 : 3 = 60^\circ$.
}
\end{ex}

\begin{ex}
Cho góc vuông xOy. Tia Oz chia góc xOy thành hai góc bằng nhau. Mỗi góc nhỏ có số đo bao nhiêu độ?
\shortans{45}
\loigiai{
$90 : 2 = 45^\circ$.
}
\end{ex}

\begin{ex}
Cho góc bẹt mOn. Tia Ot chia góc mOn thành hai góc bằng nhau. Mỗi góc nhỏ có số đo bao nhiêu độ?
\shortans{90}
\loigiai{
$180 : 2 = 90^\circ$.
}
\end{ex}

\begin{ex}
Dùng ê-ke để kiểm tra góc vuông có bao nhiêu góc vuông trên chiếc ê-ke?
\shortans{1}
\loigiai{
Một chiếc ê-ke có 1 góc vuông.
}
\end{ex}

\begin{ex}
Nếu kim phút quay từ số 12 đến số 3 thì đã quay một góc bao nhiêu độ?
\shortans{90}
\loigiai{
$90^\circ$.
}
\end{ex}

\begin{ex}
Nếu kim phút quay từ số 12 đến số 6 thì đã quay một góc bao nhiêu độ?
\shortans{180}
\loigiai{
$180^\circ$.
}
\end{ex}

\begin{ex}
Nếu kim phút quay từ số 12 đến số 4 thì đã quay một góc bao nhiêu độ?
\shortans{120}
\loigiai{
$30 \times 4 = 120^\circ$.
}
\end{ex}

\begin{ex}
Nếu kim phút quay từ số 12 đến số 2 thì đã quay một góc bao nhiêu độ?
\shortans{60}
\loigiai{
$30 \times 2 = 60^\circ$.
}
\end{ex}

\begin{ex}
Một góc có số đo $75^\circ$, cần thêm bao nhiêu độ nữa để được góc vuông?
\shortans{15}
\loigiai{
$90 - 75 = 15^\circ$.
}
\end{ex}

\begin{ex}
Một góc có số đo $115^\circ$, cần thêm bao nhiêu độ nữa để được góc bẹt?
\shortans{65}
\loigiai{
$180 - 115 = 65^\circ$.
}
\end{ex}

\begin{ex}
Góc nhọn lớn nhất có số đo độ nguyên là bao nhiêu độ?
\shortans{89}
\loigiai{
Là $89^\circ$.
}
\end{ex}
